
\documentclass{article}
%%%%%%%%%%%%%%%%%%%%%%%%%%%%%%%%%%%%%%%%%%%%%%%%%%%%%%%%%%%%%%%%%%%%%%%%%%%%%%%%%%%%%%%%%%%%%%%%%%%%%%%%%%%%%%%%%%%%%%%%%%%%%%%%%%%%%%%%%%%%%%%%%%%%%%%%%%%%%%%%%%%%%%%%%%%%%%%%%%%%%%%%%%%%%%%%%%%%%%%%%%%%%%%%%%%%%%%%%%%%%%%%%%%%%%%%%%%%%%%%%%%%%%%%%%%%
%TCIDATA{OutputFilter=LATEX.DLL}
%TCIDATA{Version=5.50.0.2960}
%TCIDATA{<META NAME="SaveForMode" CONTENT="1">}
%TCIDATA{BibliographyScheme=Manual}
%TCIDATA{Created=Monday, October 19, 2026 11:29:52}
%TCIDATA{LastRevised=Tuesday, October 20, 2026 16:42:32}
%TCIDATA{<META NAME="GraphicsSave" CONTENT="32">}
%TCIDATA{<META NAME="DocumentShell" CONTENT="Scientific Notebook\Blank Document">}
%TCIDATA{CSTFile=Math with theorems suppressed.cst}
%TCIDATA{PageSetup=72,72,72,72,0}
%TCIDATA{AllPages=
%F=36,\PARA{038<p type="texpara" tag="Body Text" >\hfill \thepage}
%}


\newtheorem{theorem}{Theorem}
\newtheorem{acknowledgement}[theorem]{Acknowledgement}
\newtheorem{algorithm}[theorem]{Algorithm}
\newtheorem{axiom}[theorem]{Axiom}
\newtheorem{case}[theorem]{Case}
\newtheorem{claim}[theorem]{Claim}
\newtheorem{conclusion}[theorem]{Conclusion}
\newtheorem{condition}[theorem]{Condition}
\newtheorem{conjecture}[theorem]{Conjecture}
\newtheorem{corollary}[theorem]{Corollary}
\newtheorem{criterion}[theorem]{Criterion}
\newtheorem{definition}[theorem]{Definition}
\newtheorem{example}[theorem]{Example}
\newtheorem{exercise}[theorem]{Exercise}
\newtheorem{lemma}[theorem]{Lemma}
\newtheorem{notation}[theorem]{Notation}
\newtheorem{problem}[theorem]{Problem}
\newtheorem{proposition}[theorem]{Proposition}
\newtheorem{remark}[theorem]{Remark}
\newtheorem{solution}[theorem]{Solution}
\newtheorem{summary}[theorem]{Summary}
\newenvironment{proof}[1][Proof]{\noindent\textbf{#1.} }{\ \rule{0.5em}{0.5em}}
\input{tcilatex}
\begin{document}


Remember we studied intensity%
\[
I=\frac{P}{A} 
\]

Photographic film and digital focal plane arrays detect the intensity of
light that hits them. We can see that the area of our image depends on our
magnification, which depends on $d_{i}$ and for our distant objects it is
proportional to $f.$ The image area is proportional to $d_{i}^{2}\approx
f^{2}.$ So we can say that the area is proportional to $f^{2}.$ Then 
\[
I\propto \frac{P}{f^{2}} 
\]%
The power entering the camera is proportional to the size of the aperture
(hole the light goes through). A bigger aperture lets in more light. A
smaller aperture lets in less light. If the camera has a circular opening,
this area is proportional to the square of the diameter of the opening, $%
D^{2}$ so%
\[
I\propto \frac{D^{2}}{f^{2}} 
\]%
This ratio is useful because it gives us a relative measure of how much
intensity we get in terms of things we can easily know about our camera.
Good cameras have changeable aperture sizes, and many lenses have changeable
focal lengths. By using the combination of these two terms, we can ensure we
will get enough light (but not too much) when we take the picture.

It would be good to give this ratio a special name. But instead, we named
the ratio 
\begin{equation}
f/\#\equiv \frac{f}{D}
\end{equation}%
It is called the $f/\#$ (pronounced f-number, or f-stop) so our intensity is
proportional to one over the f-number. 
\begin{equation}
I\propto \frac{1}{\left( f/\#\right) ^{2}}
\end{equation}%
Good cameras have adjustable lens systems marked in $f/\#^{\prime }s$.
Typical values are $f/2.8,$ $f/4,$ $f/5.6,$ $f/8,$ $f/11,$ and $f/16.$
Notice that since $I$ is proportional to $\left( 1/f/\#\right) ^{2}$ these
common $f/\#$ values give an increase of intensity of a factor of $2$ each
time you change the $f/\#$ by one marked stop. That is because, say, 
\[
\frac{f/4}{f/2.8}=\frac{4}{2.8}=1.\,\allowbreak 428\,6\allowbreak =\sqrt{2} 
\]
and the intensity goes as $1/\left( f/\#\right) ^{2}$ so the change in
intensity from $f/2.8$ to $f/4$ would be proportional to $\left( \sqrt{2}%
\right) ^{2}=2.$ It is the same between any two of our standard $f/\#$
values.%
\[
\begin{tabular}{|c|c|}
\hline
$f/\#$ change & Intensity change factor \\ \hline
$f/11$ to $f/8$ & $2$ \\ \hline
$f/8$ to $f/5.6$ & $2$ \\ \hline
$f/5.6$ to $f/4$ & $2$ \\ \hline
$f/4$ to $f/2.8$ & $2$ \\ \hline
\end{tabular}%
\]

This terminology is used for telescope design as well. The Hubble telescope
is an $f/24$ Ritchey-Chretien Cassegrainian system with a $2.4\unit{m}$
diameter aperture. The effective focal length is $57.6\unit{m}$.

\FRAME{dhF}{1.8121in}{1.1229in}{0pt}{}{}{Figure}{\special{language
"Scientific Word";type "GRAPHIC";maintain-aspect-ratio TRUE;display
"USEDEF";valid_file "T";width 1.8121in;height 1.1229in;depth
0pt;original-width 3.7118in;original-height 2.29in;cropleft "0";croptop
"1";cropright "1";cropbottom "0";tempfilename
'TN66HT00.wmf';tempfile-properties "XPR";}}

It is important to realize that electronic (and biological) sensors don't
react instantly to what we see. The intensity is 
\[
I=\frac{P}{A}=\frac{\Delta E}{\Delta tA}
\]%
so there is a time involved. The time it takes to collect enough light to
form an image on the sensor is called the \emph{exposure time}. 
\[
\Delta t=\frac{\Delta E}{IA}
\]%
So changing our $f/\#$ changes the needed exposure time by changing the
intensity. 
\[
\Delta t\propto \frac{\Delta E}{\frac{1}{\left( f/\#\right) ^{2}}A}=\frac{%
\Delta E}{A}\left( f/\#\right) ^{2}
\]%
How sensitive our camera sensor is also affects the exposure time. Modern
sensors have adjustable sensitivity. The photography world gives the three
letters ISO for the name for this detector sensitivity. There isn't a
standard for exactly what ISO setting gives what exposure. Different
manufacturers use slightly different numbers. But a change in ISO settings
usually are equivalent to one $f/\#$ change in exposure.

This is part of what a good photographer does in taking a picture. The
photographer will adjust the $f/\#$ and the exposure time and ISO to get a
photograph that is not over exposed (too light) or underexposed (to dark).

\end{document}
